\documentclass[11pt,a4paper]{article}
\usepackage[utf8]{inputenc}
\usepackage[T1]{fontenc}
\usepackage{lmodern}
\usepackage[french]{babel}
\frenchsetup{AutoSpacePunctuation=false}
\usepackage[margin=2.5cm]{geometry}
\usepackage{amsmath,amssymb,amsthm}
\usepackage{enumitem}
\usepackage[colorlinks=true,linkcolor=blue!50!black,urlcolor=blue!50!black,citecolor=blue!50!black]{hyperref}
\usepackage{xcolor}
\setlength{\parskip}{3pt}
\setlength{\emergencystretch}{2em}

\newtheorem{theorem}{Théorème}
\newtheorem{proposition}[theorem]{Proposition}
\newtheorem{lemma}[theorem]{Lemme}
\newtheorem{corollary}[theorem]{Corollaire}
\theoremstyle{remark}
\newtheorem{remark}[theorem]{Remarque}

\newcommand{\Z}{\mathbb{Z}}
\newcommand{\Q}{\mathbb{Q}}
\newcommand{\F}{\mathbb{F}}
\newcommand{\one}{\mathbf{j}}
\newcommand{\srg}{\operatorname{srg}}
\newcommand{\Fix}{\operatorname{Fix}}
\newcommand{\tr}{\operatorname{tr}}
\newcommand{\rk}{\operatorname{rang}}

\title{Le graphe de Conway $\srg(99,14,1,2)$ n'a pas d'automorphisme d'ordre~3\\[3pt]
\large Une obstruction par les discriminants}
\author{Dilan Djematene}
\date{23 septembre 2026 \\[4pt]
\textcolor{red!60!black}{\textbf{Preuve candidate, non encore relue par un tiers.}}\\[2pt]
\small Tout est vérifié en Lean~4 / Mathlib (voir §6), sauf le corollaire~8, qui s'appuie sur des résultats publiés.}

\begin{document}
\maketitle

\begin{abstract}
Soit $\Gamma$ un graphe fortement régulier de paramètres $(99,14,1,2)$ et $g$ un automorphisme d'ordre~3 sans point fixe.
L'action de $g$ fournit deux matrices $33\times 33$ : le quotient $Q$, à coefficients entiers, et une matrice hermitienne $M$ à
coefficients dans les entiers d'Eisenstein $\Z[\omega]$, congrues modulo $1-\omega$. On montre qu'un tel couple ne peut pas exister :
les réseaux des vecteurs propres de valeur propre~$3$ de $Q$ et de $M$ ont des réductions modulo~3 égales, donc des
déterminants congrus modulo~3. Or le premier vaut $2$ modulo~3, à cause du facteur $11$ de $\langle\one,\one\rangle=33$,
et le second vaut $1$, car seul le premier $7$ y intervient. Le cas d'un triangle fixe se traite de la même façon, une fois
observé que les coordonnées des points fixes sont divisibles par~3.
Un comptage élémentaire (qui redémontre un résultat de Makhnev et Minakova) montre qu'un automorphisme d'ordre~3 a au plus trois
points fixes ; il en résulte que $\Gamma$ n'a aucun automorphisme d'ordre~3.
Le cas sans point fixe était le seul cas d'ordre premier laissé ouvert par la recherche sur ordinateur de Behbahani et Lam ;
joint aux résultats publiés, on obtient que le groupe d'automorphismes de $\Gamma$ est trivial ou d'ordre~2.
La preuve est de l'algèbre linéaire élémentaire sur $\Z$ et $\Z[\omega]$, sans calcul par ordinateur ; elle est entièrement
vérifiée en Lean~4, à partir de la définition d'un graphe fortement régulier, sans aucune entrée extérieure.
\end{abstract}

%====================================================================
\section{Énoncés}

Soit $A$ la matrice d'adjacence de $\Gamma=\srg(99,14,1,2)$ ; $A^2=14I+A+2(J-I-A)$, soit
\begin{equation}\label{eq:srg}
A^2+A=12I+2J ,
\end{equation}
où $J$ est la matrice dont tous les coefficients valent~$1$. On note $\one$ le vecteur dont toutes les coordonnées valent~$1$,
$\omega=e^{2i\pi/3}$ et $\pi=1-\omega$.

\begin{theorem}\label{thm:libre}
$\Gamma$ n'a aucun automorphisme d'ordre~$3$ sans point fixe.
\end{theorem}

\begin{proposition}\label{prop:triangle}
$\Gamma$ n'a aucun automorphisme d'ordre~$3$ ayant exactement trois points fixes (ils formeraient un triangle).
\end{proposition}

\begin{corollary}\label{cor:tout}
$\Gamma$ n'a aucun automorphisme d'ordre~$3$.
\end{corollary}
\begin{proof}
Soit $g\ne1$ avec $g^3=1$. Par le lemme~\ref{lem:fix} (§5), $g$ a au plus trois points fixes. Leur nombre est
$\equiv 99\equiv0\pmod3$, car les autres sommets forment des orbites de taille~3 : il vaut donc $0$ ou $3$, et l'on conclut par le
théorème~\ref{thm:libre} et la proposition~\ref{prop:triangle}. (Le lemme~\ref{lem:fix} est dû à Makhnev et Minakova \cite{MM04},
sous la forme « les points fixes forment l'ensemble vide ou un triangle » ; le cas avec points fixes est aussi exclu, par une
recherche sur ordinateur, dans \cite[th.~4.14]{B09} et \cite{BL11}.)
\end{proof}

\paragraph{Contexte.} L'existence d'un $\srg(99,14,1,2)$ est un problème ouvert ; c'est l'un des cinq problèmes à 1\,000 dollars de
Conway \cite{Co}. Voir \cite{W84} pour une première étude de ce graphe et \cite{BVM22} pour le cadre général. Sur le groupe
$G=\mathrm{Aut}(\Gamma)$, on savait ceci.
\begin{itemize}[leftmargin=*,nosep]
\item Makhnev et Minakova \cite{MM04}, par la théorie des caractères : $|G|$ divise $2\cdot3^3\cdot7\cdot11$, et un
automorphisme d'ordre~3 fixe soit aucun sommet, soit les trois sommets d'un triangle.
\item Behbahani et Lam \cite{B09,BL11}, par une recherche sur ordinateur (matrices d'orbites) : seuls $2$ et $3$ peuvent diviser
$|G|$, et un automorphisme d'ordre~3 n'a aucun point fixe \cite[th.~4.14]{B09}. Le cas d'un automorphisme d'ordre~3 sans point
fixe est le seul que ce calcul n'a pas tranché \cite[table~21]{B09}.
\item Crnković et Maksimović \cite{CM20}, par ordinateur : $\Gamma$ n'admet aucun groupe d'automorphismes d'ordre $6$ ou $9$.
\item Makhnev \cite[cor.~2.4]{Mak10} : si $2$ divise $|G|$, alors $|G|$ divise $42$. Cesarz et Woldar \cite{CW25}, sans
ordinateur : si $7$ divise $|G|$, alors $G\simeq\Z/7\Z$ ; par suite, si $2$ divise $|G|$, alors $|G|$ divise $6$.
\item Ishida \cite[prop.~8.7]{I26}, par une identité trace--rang : un automorphisme d'ordre~3 sans point fixe a exactement six
orbites qui induisent un triangle.
\end{itemize}
Réunis, ces résultats donnent $|G|\in\{1,2,3\}$ : le seul cas encore ouvert était un automorphisme d'ordre~3 sans point fixe.
Il est encore déclaré non résolu dans \cite{TS26}, après une tentative par solveur CP-SAT (voir aussi \cite{K26} pour une
approche par solveur SAT sans hypothèse de symétrie). Le théorème~\ref{thm:libre} le ferme ; la proposition~\ref{prop:triangle}
et le lemme~\ref{lem:fix} redémontrent, sans ordinateur, le cas avec points fixes. Le corollaire~\ref{cor:aut} (§5) en tire les
conséquences pour $|G|$. L'argument exploite l'invariant suivant, le \emph{discriminant}, et retrouve au passage le nombre six
d'Ishida (remarque~\ref{rem:tau}).

%====================================================================
\section{Réduction à un énoncé sur des matrices}

Soit $g$ un automorphisme d'ordre~3 sans point fixe. Ses $33$ orbites ont chacune trois éléments ; on choisit un représentant dans
chacune et l'on note les sommets $(U,b)$, $U\in\{1,\dots,33\}$, $b\in\Z/3$, avec $g(U,b)=(U,b+1)$. Pour $b\in\Z/3$, soit
$A_b\in M_{33}(\Z)$ définie par $A_b[U,V]=A[(U,0),(V,b)]$. Comme $g$ est un automorphisme, $A[(U,c),(V,c+b)]=A_b[U,V]$ pour tout
$c$, et la symétrie de $A$ donne $A_b^{\mathsf T}=A_{-b}$. On pose
\[
Q=A_0+A_1+A_2\in M_{33}(\Z),\qquad M=A_0+\omega A_1+\omega^2A_2\in M_{33}(\Z[\omega]).
\]

\begin{lemma}\label{lem:QM}
(i) $Q$ est symétrique, $Q\one=14\,\one$ et $Q^2+Q=12I+6J$.
(ii) $M$ est hermitienne et $M^2+M=12I$.
(iii) $Q-M=(1-\omega)A_1+(1-\omega^2)A_2$ est divisible par $\pi$ coefficient par coefficient.
(iv) $\tr Q=2\tau$ et $\tr M=-\tau$, où $\tau$ est le nombre d'orbites qui induisent un triangle.
\end{lemma}
\begin{proof}
Pour une matrice $X$ indexée par les sommets et commutant avec $g$, posons $X_b[U,V]=X[(U,0),(V,b)]$ et
$\Phi(X)=\sum_b X_b\,x^b$, élément de $M_{33}\bigl(\Z[x]/(x^3-1)\bigr)$. On a ${(XY)_b=\sum_c X_cY_{b-c}}$, donc $\Phi$ est un
morphisme d'anneaux,
avec $\Phi(I)=I$ et $\Phi(J)=(1+x+x^2)J$. Appliquer $\Phi$ à \eqref{eq:srg}, puis évaluer en $x=1$ et en $x=\omega$
($1+\omega+\omega^2=0$), donne les équations de (i) et (ii). Les sommes de lignes de $Q$ sont les degrés, d'où $Q\one=14\one$.
On a $Q^{\mathsf T}=\sum_b A_{-b}=Q$ et $M^*=\sum_b\bar\omega^{\,b}A_{-b}=M$. Le point (iii) vient de $1-\omega^2=(1-\omega)(1+\omega)$.
Pour (iv) : $(U,0)$ est voisin de $(U,1)$ si et seulement s'il l'est de $(U,2)$, et c'est le cas exactement pour les orbites
triangulaires ; d'où $Q[U,U]=2$ ou $0$ et $M[U,U]=\omega+\omega^2=-1$ ou $0$.
\end{proof}

Le théorème~\ref{thm:libre} découle donc du lemme~\ref{lem:QM} et de l'énoncé suivant.

\begin{theorem}\label{thm:matrices}
Il n'existe pas de couple $(Q,M)$ tel que
$Q\in M_{33}(\Z)$ soit symétrique avec $Q\one=14\one$ et $Q^2+Q=12I+6J$, $M\in M_{33}(\Z[\omega])$ soit hermitienne avec $M^2+M=12I$,
et $Q\equiv M\pmod{\pi}$ coefficient par coefficient.
\end{theorem}

%====================================================================
\section{Preuve du théorème~\ref{thm:matrices}}

Soit $\rho:\Z[\omega]\to\F_3$ le morphisme de réduction ($\rho(\omega)=1$, de noyau $\pi\Z[\omega]$). Sur $\Z$, c'est la
réduction modulo~3, et $\rho(\bar z)=\rho(z)$. Pour une matrice $X$, on note $\bar X=\rho(X)$ ; alors $\rho(X^*)=\bar X^{\mathsf T}$.

\paragraph{Étape 1 (identités).}
Soit $V=\ker(Q-3I)\subset\Q^{33}$.
(a) Si $x\in V$, alors $\one^{\mathsf T}x=0$ : en effet $14\,\one^{\mathsf T}x=\one^{\mathsf T}Qx=3\,\one^{\mathsf T}x$.
(b) La matrice entière symétrique $N=11(Q+4I)-6J$ vérifie ${(Q-3I)N=0}$ et $Nx=77x$ pour $x\in V$.
En effet, $(Q-3I)(Q+4I)=Q^2+Q-12I=6J$ et $(Q-3I)J=11J$, donc $(Q-3I)N=66J-66J=0$ ; si $x\in V$, $(Q+4I)x=7x$ et $Jx=0$.
(c) La matrice hermitienne $N'=M+4I$ vérifie $(M-3I)N'=M^2+M-12I=0$ et $N'y=7y$ pour $y\in\ker(M-3I)$.
(d) Comme $Q$ est symétrique avec $Q\one=14\one$, on a $QJ=JQ=14J$ ; sur $\one^\perp$, $Q^2+Q=12I$, donc $14$ est valeur propre
simple et $W:=\ker\bigl((Q-3I)(Q-14I)\bigr)=V\oplus\Q\one$. La matrice entière symétrique $N_W=3(Q+4I)-J$ vérifie
$(Q-3I)(Q-14I)N_W=18(Q-14I)J=0$ et $N_Ww=21w$ pour $w\in W$ (sur $V$ par (a) ; et $N_W\one=54\one-33\one$).

\paragraph{Étape 2 (bases entières).}
Soit $R$ un anneau principal ($\Z$ ou $\Z[\omega]$) et $X\in M_n(R)$. Le module $L=\{v\in R^n : Xv=0\}$ est un facteur direct de
$R^n$ : en effet $R^n/L\simeq X(R^n)\subset R^n$ est de type fini sans torsion, donc libre. Il existe donc une matrice $B$ de taille
$n\times r$ dont les colonnes forment une base de $L$ ($r$ est la dimension du noyau sur le corps des fractions), et une matrice
$C$ de taille $r\times n$ telle que $CB=I_r$ ; alors $BCv=v$ pour tout $v\in L$. On applique ceci :
\begin{itemize}[nosep]
\item à $L=\ker_\Z(Q-3I)$ : $(B,C,r)$, $G=B^{\mathsf T}B$, $d=\det G>0$ ;
\item à $L'=\ker_{\Z[\omega]}(M-3I)$ : $(B',C',r')$ et $G'=B'^*B'$, matrice hermitienne définie positive de déterminant
${d'=\det G'>0}$ ;
comme $d'$ est réel et appartient à $\Z[\omega]$, $d'\in\Z$ ;
\item à $L_W=\ker_\Z\bigl((Q-3I)(Q-14I)\bigr)$, de rang $r+1$ : $(B_W,C_W)$, $G_W=B_W^{\mathsf T}B_W$.
\end{itemize}

\paragraph{Étape 3 (les identités $GK=77I$ et $G'K'=7I$).}
Posons $K=CNC^{\mathsf T}\in M_r(\Z)$. Les colonnes de $N$ sont dans $L$ (étape 1b), donc $BCN=N$ ; comme $N$ est symétrique,
$N=NC^{\mathsf T}B^{\mathsf T}$, d'où $N=BKB^{\mathsf T}$. De plus $NB=77B$, donc $B^{\mathsf T}N=77B^{\mathsf T}$ et
\[
GK=B^{\mathsf T}(BCN)C^{\mathsf T}=B^{\mathsf T}NC^{\mathsf T}=77\,(CB)^{\mathsf T}=77\,I_r .
\]
De même, avec $K'=C'N'C'^*$ et $K_W=C_WN_WC_W^{\mathsf T}$ : $N'=B'K'B'^*$, $G'K'=7I_{r'}$ et $G_WK_W=21I_{r+1}$. Conséquences :
\begin{itemize}[nosep]
\item $d\cdot\det K=77^r$, donc $d=7^a\,11^b$ ;
\item $K'$ est hermitienne, donc $\det K'\in\Z$, et $d'\cdot\det K'=7^{r'}$ ; donc $d'=7^c$, et $d'\equiv 1\pmod 3$ ;
\item $\det G_W$ divise $21^{r+1}$ ; en particulier $11\nmid\det G_W$.
\end{itemize}

\paragraph{Étape 4 (l'exposant $b$ est impair).}
Les colonnes de la matrice $[B\,|\,\one]$ sont dans $L_W$, donc $[B\,|\,\one]=B_WT$ avec $T=C_W[B\,|\,\one]\in M_{r+1}(\Z)$.
Par l'étape 1(a), $B^{\mathsf T}\one=0$ : la matrice de Gram de $[B\,|\,\one]$ vaut $\mathrm{diag}(G,33)$, et aussi
$T^{\mathsf T}G_WT$. Donc
\[
33\,d=\det(T)^2\,\det G_W .
\]
Les valuations $11$-adiques donnent $1+b=2\,v_{11}(\det T)$ : $b$ est impair. Donc $d=7^a11^b\equiv(-1)^b=-1\pmod 3$.

\paragraph{Étape 5 (réduction modulo~3 : $d\equiv d'$).}
Comme $Q\equiv M\pmod\pi$, $\bar N=2(\bar Q+I)$ et $\bar N'=\bar Q+I$ ; donc $\bar N=2\bar N'$ et $\bar N'=2\bar N$. En réduisant les
identités de l'étape~3 :
\[
\bar K=\bar C\bar N\bar C^{\mathsf T}=2\,(\bar C\bar B')\,\bar K'\,(\bar C\bar B')^{\mathsf T},\qquad
\bar K'=\bar C'\bar N'\bar C'^{\mathsf T}=2\,(\bar C'\bar B)\,\bar K\,(\bar C'\bar B)^{\mathsf T}.
\]
Les déterminants $\det K=77^r/d$ et $\det K'=7^{r'}/d'$ sont premiers à~3, donc $\bar K$ et $\bar K'$ sont inversibles. La première
identité factorise $\bar K$ par $\F_3^{\,r'}$, donc $r\le r'$ ; la seconde donne $r'\le r$. Ainsi $r=r'$, $\bar C\bar B'$ est carrée
et inversible, et en prenant les déterminants dans la première identité :
\[
2^r\,\rho(d)^{-1}=\det\bar K=2^r\,\det(\bar C\bar B')^2\,\det\bar K'=2^r\,\rho(d')^{-1},
\]
car un carré non nul de $\F_3$ vaut $1$, $\rho(77)=2$ et $\rho(7)=1$. Donc $d\equiv d'\pmod 3$.

\paragraph{Étape 6.} Les étapes 4 et 3 donnent $d\equiv 2$ et $d'\equiv 1\pmod 3$, ce qui contredit l'étape~5. \qed

\begin{remark}\label{rem:tau}
$\tr N=\tr(BKB^{\mathsf T})=\tr(KG)=77r$ et de même $\tr N'=7r'$. Avec le lemme~\ref{lem:QM}(iv), $77r=11(2\tau+132)-198$ et
$7r'=132-\tau$. L'égalité $r=r'$ obtenue à l'étape~5 donne donc $2\tau+114=132-\tau$, soit $\tau=6$ et $r=r'=18$ : c'est le
résultat d'Ishida, retrouvé par la même réduction. L'obstruction nouvelle est l'étape suivante : l'égalité des
\emph{discriminants}.
\end{remark}

%====================================================================
\section{Le cas de trois points fixes (proposition~\ref{prop:triangle})}

Soit $g$ d'ordre~3 ayant exactement trois points fixes $t_1,t_2,t_3$. Ils forment un triangle : si $t_i$ et $t_j$ n'étaient pas
adjacents, $g$ permuterait leurs $\mu=2$ voisins communs, donc les fixerait (il est d'ordre~3), et aurait au moins quatre points
fixes. Les $96$ autres sommets forment $32$ orbites libres, de sommets $(U,b)$, $U\le32$ ; on définit $A_b$, $Q$ et $M$
($32\times32$) comme au §2, et $\mathbf a[U,i]=A[(U,0),t_i]$. Alors $M$ est hermitienne, $M^2+M=12I$ (dans le calcul du §2, les
chemins passant par un $t_i$ et le terme en $J$ sont multipliés par $1+\omega+\omega^2=0$) et $Q\equiv M\pmod\pi$.

Les vecteurs de $\Z^{99}$ invariants par $g$ s'identifient à $\Z^{35}$ (une valeur par orbite libre, une par $t_i$). $A$ y agit par
\[
\tilde A=\begin{pmatrix}Q&\mathbf a\\3\mathbf a^{\mathsf T}&J_3-I_3\end{pmatrix},\qquad
\tilde A\one=14\one,\qquad \tilde A^2+\tilde A=12I+2\,\one w^{\mathsf T},
\]
où $w=(3,\dots,3,1,1,1)$ ($J$ agit par $\one w^{\mathsf T}$). Le produit scalaire de $\Z^{99}$ induit la forme $D=\mathrm{diag}(w)$,
pour laquelle $\tilde A$ est autoadjointe ($D\tilde A$ est symétrique). On pose $\Delta=\mathrm{diag}(I_{32},3I_3)$, de sorte que
$D\Delta=3I$.

\begin{enumerate}[leftmargin=*,nosep]
\item \emph{Les coordonnées fixes sont divisibles par~3.} Soit $\tilde N=11(\tilde A+4I)-2\,\one w^{\mathsf T}$. Comme à
l'étape~1, $(\tilde A-3I)\tilde N=0$ et $\tilde Nx=77x$ si $\tilde Ax=3x$ (alors $w^{\mathsf T}x=0$). Les trois dernières lignes
de $\tilde N$ sont nulles modulo~3 : $\tilde N[t_i,U]=33\,\mathbf a[U,i]-6$ et $\tilde N[t_i,t_j]=9+33\,\delta_{ij}$. Si
$x\in L=\ker_\Z(\tilde A-3I)$, on a donc $77\,x_{t_i}\equiv0$, soit $3\mid x_{t_i}$.
\item \emph{Gram pondéré.} Soit $B$ une base de $L$ (rang $r$) et $C$ telle que $CB=I$. Par le point~1, $\tilde G=\tfrac13
B^{\mathsf T}DB$ est une matrice entière, définie positive. La matrice $X=\tilde N\Delta$ est symétrique, $DX=3\tilde N^{\mathsf T}$,
et ses colonnes sont dans $L$. Comme à l'étape~3, $X=BKB^{\mathsf T}$ avec $K=CXC^{\mathsf T}$ entière, et
$3\tilde GK=B^{\mathsf T}DXC^{\mathsf T}=3(\tilde NB)^{\mathsf T}C^{\mathsf T}=3\cdot77\,I$. Donc $\tilde GK=77I$ et
$\det\tilde G=7^a11^b$.
\item \emph{Réduction modulo~3.} On a $X\equiv\mathrm{diag}\bigl(2(\bar Q+I),0\bigr)$ (les trois dernières colonnes sont
multipliées par~3, les trois dernières lignes sont nulles) et $\bar B=\binom{\bar B_1}{0}$ par le point~1. Réduites modulo~3, les
identités $X=BKB^{\mathsf T}$ et $K=CXC^{\mathsf T}$ ne font donc intervenir que le bloc $32\times32$ :
$2(\bar Q+I)=\bar B_1\bar K\bar B_1^{\mathsf T}$ et $\bar K=2\,\bar C_1(\bar Q+I)\bar C_1^{\mathsf T}$. Du côté de $M$, les
étapes~2 et~3 s'appliquent sans changement ($d'=7^c\equiv1$), et l'étape~5, mot pour mot, donne $r=r'$ et
$\det\tilde G\equiv d'\equiv1\pmod3$.
\item \emph{Le premier 11.} $W=\ker_\Z\bigl((\tilde A-3I)(\tilde A-14I)\bigr)$ contient $L$ et $\one$ et est de rang $r+1$ (sur
$\{w^{\mathsf T}x=0\}$, $\tilde A^2+\tilde A=12I$). La matrice $N_W=9(\tilde A+4I)-\one w^{\mathsf T}$ vaut $63$ fois le projecteur
sur $W\otimes\Q$ ; avec $X_W=N_W\Delta$, le calcul du point~2 donne $G_WK_W=189\,I$ pour ${G_W=B_W^{\mathsf T}DB_W}$, donc
$11\nmid\det G_W$. Comme $w^{\mathsf T}B=0$, la matrice de Gram (pour $D$) de $[B\,|\,\one]$ est
$\mathrm{diag}(3\tilde G,\,\one^{\mathsf T}D\one)=\mathrm{diag}(3\tilde G,99)$ ; comme à l'étape~4,
$3^r\cdot99\cdot\det\tilde G=\det(T)^2\det G_W$. Donc $b$ est impair et $\det\tilde G\equiv2\pmod3$, ce qui contredit le
point~3. \qed
\end{enumerate}

%====================================================================
\section{Les points fixes (corollaire~\ref{cor:tout}) et l'ordre de $\mathrm{Aut}(\Gamma)$}

\begin{lemma}\label{lem:fix}
Soit $g\ne1$ un automorphisme de $\Gamma$ tel que $g^3=1$. Alors $g$ a au plus trois points fixes.
\end{lemma}
\begin{proof}
Soit $F$ l'ensemble des points fixes et, pour $x\in F$, $n_x$ le nombre de voisins fixes de $x$.
\begin{enumerate}[leftmargin=*,nosep]
\item \emph{Fermeture.} Un sommet $v$ adjacent à deux points fixes distincts $a,b$ est fixe : si $a\sim b$, $v$ est leur unique
voisin commun ($\lambda=1$) ; sinon $g$ permute les $\mu=2$ voisins communs de $a$ et $b$ et, étant d'ordre~3, les fixe. En
particulier, un sommet non fixe a au plus un voisin fixe.
\item \emph{$n_x\in\{2,8,14\}$.} Les voisins non fixes de $x\in F$ forment des orbites de taille~3, donc $n_x\equiv14\equiv2\pmod3$.
Et $n_x$ est pair : par $\lambda=1$, un voisin fixe $u$ de $x$ a un unique voisin commun avec $x$, qui est fixe par le point~1 ;
cela apparie les voisins fixes de $x$.
\item \emph{Comptage.} Par le point~1, il y a au plus $99-|F|$ arêtes entre $F$ et son complémentaire :
$\sum_{x\in F}(14-n_x)\le 99-|F|$, soit $15\,|F|\le 99+\sum_{x\in F}n_x$.
\item \emph{Conclusion.} Si $n_x=14$, tout sommet non adjacent à $x$ a deux voisins communs avec $x$, fixes, donc est fixe
(point~1) : $g=1$, exclu. Si $n_x=8$ et qu'aucun $n_y$ ne vaut $14$, soit $A$ l'ensemble des huit voisins fixes de $x$ et
$S=F\setminus(\{x\}\cup A)$, de cardinal $s$. Un $u\in A$ a pour voisins fixes $x$, au plus un sommet de $A$ ($\lambda=1$) et des
sommets de $S$ ; un $y\in S$ a au plus $\mu=2$ voisins dans $A$. Donc $\sum_{u\in A}n_u\le16+2s$ et
$\sum_{y\in F}n_y\le 8+(16+2s)+8s$. Comme $|F|=9+s$, le point~3 donne $15(9+s)\le 123+10s$, soit $12+5s\le0$ : impossible. Reste
le cas où tous les $n_y$ valent~$2$. Si $y\in F$ n'est ni $x$ ni un voisin de $x$, un voisin commun $a$ de $x$ et $y$ est fixe,
ainsi que le voisin commun $c$ de $x$ et $a$ (point~1) ; alors $a$ a trois voisins fixes $x,y,c$, contradiction. Donc tout point
fixe est $x$ ou l'un de ses deux voisins fixes : $|F|\le3$. \qedhere
\end{enumerate}
\end{proof}

Joint aux résultats publiés rappelés au §1, le corollaire~\ref{cor:tout} donne ceci.

\begin{corollary}\label{cor:aut}
Supposons que $\Gamma$ existe, et soit $G=\mathrm{Aut}(\Gamma)$.
\begin{enumerate}[label=(\roman*),nosep]
\item Avec \cite{MM04} et \cite{CW25}, obtenus sans ordinateur : $|G|\in\{1,2,7,11\}$.
\item Avec en plus la recherche sur ordinateur de \cite{B09,BL11} : $|G|\in\{1,2\}$, c'est-à-dire que $G$ est trivial ou
isomorphe à $\Z/2\Z$.
\end{enumerate}
\end{corollary}
\begin{proof}
Par le corollaire~\ref{cor:tout}, $G$ n'a aucun élément d'ordre~3, donc $3\nmid|G|$ (théorème de Cauchy).
(i) Par \cite{MM04}, $|G|$ divise alors $2\cdot7\cdot11$. Si $7$ divise $|G|$, alors $G\simeq\Z/7\Z$ \cite{CW25}. Si $2$ divise
$|G|$, alors $|G|$ divise $6$ \cite{CW25}, donc $|G|=2$. Sinon, $|G|$ divise $11$.
(ii) Par \cite{BL11}, $|G|$ est alors une puissance de~$2$ ; si elle n'est pas~$1$, elle divise~$6$ \cite{CW25}, donc vaut~$2$.
\end{proof}

Un rapport du programme MIT PRIMES \cite{LM14}, dont nous n'avons pas vérifié les preuves, exclut aussi l'ordre~$11$ par un
argument élémentaire ; le point~(i) deviendrait alors $|G|\in\{1,2,7\}$. Le corollaire~\ref{cor:aut} n'est pas formalisé :
les résultats de \cite{MM04,B09,BL11,CW25} sont admis.

%====================================================================
\section{Remarques}

\begin{itemize}[leftmargin=*]
\item \textbf{Mécanisme.} Les deux caractères de $C_3$ voient la même forme quadratique $(\bar Q+I)$ sur $\F_3$. Seul le
caractère trivial contient le vecteur $\one$ ; sa norme $33=3\cdot11$ injecte le premier $11$, qui n'est pas un carré modulo~3.
Rien ne peut compenser : le seul autre premier en jeu est $7=r-s$, qui vaut $1$ modulo~3 (il se décompose dans $\Z[\omega]$).
\item \textbf{Ce qui est utilisé.} Pour le théorème~\ref{thm:libre} : uniquement \eqref{eq:srg} et l'action libre ; ni les valeurs
$0,1,2$ des coefficients de $Q$, ni aucun calcul. Pour la proposition~\ref{prop:triangle} : de même (plus $\mu=2$ pour le
triangle). Une première version utilisait l'identité trace--rang d'Ishida pour obtenir $r=r'$ ; la divisibilité par~3 des
coordonnées fixes (point~1 du §4) la rend inutile. Pour le corollaire : en plus, le comptage du lemme~\ref{lem:fix}
($\lambda=1$, $\mu=2$, $v=99$) ; aucune entrée extérieure.
\item \textbf{Généralisation.} La même preuve donne un critère pour tout $\srg(v,k,\lambda,\mu)$ de valeurs propres $k,\theta,\theta'$,
lorsque $3\nmid\theta-\theta'$ et que tous les premiers divisant $\theta-\theta'$ valent $1$ modulo~3 (hypothèses d'intégralité
vérifiées au cas par cas). Pour $v\le1300$, elle exclut
aussi les automorphismes d'ordre~3 sans point fixe pour $(105,52,21,30)$, $(105,52,29,22)$, $(261,52,11,10)$, $(714,248,97,80)$,
$(1221,610,279,330)$ et $(1221,610,329,280)$. Tous ces jeux sont d'existence inconnue dans les tables de Brouwer \cite{Br}.
\item \textbf{Contrôles numériques.} Le lemme « $d\equiv d'\pmod3$ » a été testé sans échec sur les graphes triangulaires
$T(n)$, $n\in\{6,7,9,10,12,13,15,18,21\}$, munis d'une action libre d'ordre~3 (un cas réel donne $d\equiv d'\equiv2$ : $T(10)$),
et sur environ 1\,650 réseaux invariants aléatoires, avec un recalcul indépendant ; le test échoue bien lorsque le projecteur
n'est pas 3-entier. La forme « $d=$ (premiers de $\theta-\theta'$) $\times$ (premiers de $\langle\one,\one\rangle$) » est
observée exactement sur $T(15)$, $T(18)$ et $T(21)$.
\item \textbf{Vérification formelle.} Tous les énoncés, sauf le corollaire~\ref{cor:aut}, sont vérifiés en Lean~4 \cite{dMU21}
avec Mathlib \cite{Mat20} (version \texttt{v4.35.0-rc2}), sans \texttt{sorry}, à partir de la définition de Mathlib d'un graphe
fortement régulier (\texttt{SimpleGraph.\allowbreak IsSRGWith}) :
\begin{itemize}[leftmargin=1.2em,nosep]
\item \texttt{C3.no\_order3} : tout automorphisme $\sigma$ tel que $\sigma^3=1$ est l'identité (corollaire~\ref{cor:tout},
sans \cite{MM04}) ;
\item \texttt{C3.no\_free\_order3} (théorème~\ref{thm:libre}) ;
\item \texttt{C3.no\_triangle\_order3} (proposition~\ref{prop:triangle} ; le triangle y est démontré) ;
\item \texttt{C3.no\_order3\_fix\_le3} (au plus trois points fixes) ;
\item \texttt{C3.no\_pair} (théorème~\ref{thm:matrices}) et \texttt{C3.Tri.no\_triangle\_pair} (§4), les versions matricielles.
\end{itemize}
Tous ne dépendent que des axiomes standard \texttt{propext}, \texttt{Classical.choice} et \texttt{Quot.sound}.
L'énoncé final est isolé dans le fichier \texttt{c3formal/\allowbreak Challenge.lean}, qui n'importe que Mathlib et ne contient
aucune définition. Il y est écrit sous quatre formes : graphe ; aucun élément d'ordre~3 dans $\mathrm{Aut}(\Gamma)$ ; $3$ ne
divise pas $|\mathrm{Aut}(\Gamma)|$ ; matrice d'adjacence booléenne sur $\{0,\dots,98\}$, sans bibliothèque de graphes. Trois
tests de non-vacuité l'accompagnent : les mêmes formules, avec $(9,4,1,2)$ au lieu de $(99,14,1,2)$, sont démontrées fausses
(le graphe de la tour $3\times3$ est un $\srg(9,4,1,2)$ qui a un automorphisme d'ordre~3). L'outil \texttt{lake comparator} de
Lean~4.35 contrôle que les preuves (\texttt{c3formal/\allowbreak Solution.lean}, déduites de \texttt{C3.no\_order3})
démontrent exactement ces énoncés, avec ces seuls axiomes, et qu'elles sont acceptées, avec toutes leurs dépendances jusqu'au
c\oe ur de Lean (plus de $26\,500$ déclarations), par le noyau de Lean et par cinq autres vérificateurs livrés avec Lean : \texttt{leanchecker-paranoid},
\texttt{lean4lean} \cite{Car24}, \texttt{nanoda} \cite{Nan}, \texttt{con-leche} et \texttt{con-ron}. Verdict : accepté, sur la
machine d'origine et, dans le bac à sable de l'outil, sur une machine Linux neuve (GitHub Actions). Un relecteur n'a donc à
contrôler que \texttt{c3formal/\allowbreak Challenge.lean}. Sources : dossier \texttt{c3formal/} ; détails et journaux :
\texttt{AUDIT\_ENONCE.md}.
\item \textbf{À relire en priorité (preuve papier).} L'étape~2 (facteur direct sur $\Z[\omega]$), l'étape~4 (valuation 11-adique),
l'étape~5 (réduction), au §4 le point~1, et le lemme~\ref{lem:fix}. Le corollaire~\ref{cor:aut} repose en outre sur
\cite{MM04,B09,BL11,CW25}, que nous n'avons pas revérifiés.
\end{itemize}

{\small\raggedright
\section*{Utilisation de systèmes d'IA}
Ce travail a été réalisé avec l'assistance d'un système d'IA (Claude, Anthropic), utilisé pour développer les
arguments, écrire et vérifier les programmes et la formalisation Lean~4, et rédiger le texte. L'auteur a dirigé le
travail et en assume l'entière responsabilité.

\begin{thebibliography}{BVM22}
\bibitem[B09]{B09} M.~Behbahani, \emph{On strongly regular graphs}, thèse de doctorat, Concordia University, Montréal, 2009.
\bibitem[BL11]{BL11} M.~Behbahani, C.~Lam, \emph{Strongly regular graphs with non-trivial automorphisms}, Discrete Math.\
\textbf{311} (2011), n\textsuperscript{o}~2--3, 132--144,
\href{https://doi.org/10.1016/j.disc.2010.10.005}{doi:10.1016/j.disc.2010.10.005}.
\bibitem[Br]{Br} A.\,E.~Brouwer, \emph{Parameters of strongly regular graphs}, tables en ligne,
\url{https://aeb.win.tue.nl/graphs/srg/srgtab.html}, consultées en septembre 2026.
\bibitem[BVM22]{BVM22} A.\,E.~Brouwer, H.~Van~Maldeghem, \emph{Strongly Regular Graphs}, Encyclopedia of Mathematics and its
Applications~\textbf{182}, Cambridge University Press, 2022,
\href{https://doi.org/10.1017/9781009057226}{doi:10.1017/9781009057226}.
\bibitem[Car24]{Car24} M.~Carneiro, \emph{Lean4Lean: Verifying a typechecker for Lean, in Lean}, arXiv:2403.14064, 2024
(version~3, 2025).
\bibitem[CM20]{CM20} D.~Crnković, M.~Maksimović, \emph{Construction of strongly regular graphs having an automorphism group of
composite order}, Contrib.\ Discrete Math.\ \textbf{15} (2020), n\textsuperscript{o}~1, 22--41,
\href{https://doi.org/10.55016/ojs/cdm.v15i1.62323}{doi:10.55016/ojs/cdm.v15i1.62323}.
\bibitem[Co]{Co} J.\,H.~Conway, \emph{Five \$1,000 problems (update 2017)}, OEIS Foundation,
\url{https://oeis.org/A248380/a248380.pdf}.
\bibitem[CW25]{CW25} P.\,G.~Cesarz, A.\,J.~Woldar, \emph{On the automorphism group of a putative Conway 99-graph}, Algebr.\
Comb.\ \textbf{8} (2025), n\textsuperscript{o}~2, 379--398, \href{https://doi.org/10.5802/alco.418}{doi:10.5802/alco.418} ;
arXiv:2308.02978.
\bibitem[dMU21]{dMU21} L.~de~Moura, S.~Ullrich, \emph{The Lean~4 theorem prover and programming language}, dans
\emph{Automated Deduction -- CADE~28}, Lecture Notes in Comput.\ Sci., Springer, 2021, 625--635,
\href{https://doi.org/10.1007/978-3-030-79876-5_37}{doi:10.1007/978-3-030-79876-5\_37}.
\bibitem[I26]{I26} Y.~Ishida, \emph{No involutions in the missing Moore graph}, arXiv:2606.29183v2, juillet 2026.
\bibitem[K26]{K26} A.~Keramatipour, \emph{Approaching the Conway-99 problem using SAT solvers}, arXiv:2604.23037v2, avril 2026.
\bibitem[LM14]{LM14} S.~Lou, M.~Murin, \emph{On the strongly regular graph of parameters $(99,14,1,2)$}, rapport MIT PRIMES,
2014, \url{https://math.mit.edu/research/highschool/primes/materials/2014/Lou-Murin.pdf}.
\bibitem[Mak10]{Mak10} A.\,A.~Makhnev, \emph{On automorphisms of distance-regular graphs}, J.\ Math.\ Sci.\ \textbf{166} (2010),
n\textsuperscript{o}~6, 733--742.
\bibitem[Mat20]{Mat20} The mathlib Community, \emph{The Lean mathematical library}, dans \emph{Proceedings of the 9th ACM SIGPLAN
International Conference on Certified Programs and Proofs (CPP~2020)}, ACM, 2020, 367--381,
\href{https://doi.org/10.1145/3372885.3373824}{doi:10.1145/3372885.3373824}.
\bibitem[MM04]{MM04} A.\,A.~Makhnev, I.\,M.~Minakova, \emph{On automorphisms of strongly regular graphs with parameters
$\lambda=1$, $\mu=2$}, Diskret.\ Mat.\ \textbf{16} (2004), n\textsuperscript{o}~1, 95--104,
\href{https://doi.org/10.4213/dm144}{doi:10.4213/dm144} ; trad.\ anglaise : Discrete Math.\ Appl.\ \textbf{14} (2004),
n\textsuperscript{o}~2, 201--210, \href{https://doi.org/10.1515/156939204872374}{doi:10.1515/156939204872374}.
\bibitem[Nan]{Nan} \emph{nanoda}, vérificateur de types externe pour Lean~4, \url{https://github.com/ammkrn/nanoda_lib}.
\bibitem[TS26]{TS26} A.~Thakkar, S.~Severini, \emph{A forced-structure reduction and verifiable bounds for Conway's 99-graph},
arXiv:2608.11211v2, 18~septembre 2026.
\bibitem[W84]{W84} H.\,A.~Wilbrink, \emph{On the $(99,14,1,2)$ strongly regular graph}, dans \emph{Papers dedicated to
J.\,J.~Seidel} (P.\,J.~de~Doelder, J.~de~Graaf, J.\,H.~van~Lint, éd.), EUT Report 84-WSK-03, Eindhoven University of
Technology, 1984, 342--355.
\end{thebibliography}
}

\end{document}
